\section{Proofs for Section~\ref{sec:plt2}}\label{app:plt2}

\subsection{The exponents \texorpdfstring{$1\le p<2$}{1 <= p < 2}}

\begin{proof}[Proof of \cref{prop:low}]
At $t=1$, \cref{lem:ends} gives $h(s)\ge c_0+ms$ with $m=\tfrac p2h(1)\ge0$ and
$c_0=(1-\tfrac p2)h(1)$, so that $(2-p)\abs m=pc_0$; there is equality at $s=1$. The point
$t=-1$ is the mirror image.

Let $\abs t<1$, $H=h(t)\ge0$, and let $m$ be the slope of a supporting line of $\tilde h$ at
$t$: $h(s)\ge H+m(s-t)$ on $[-1,1]$. With $c_0=H-mt$ we have to show that
\[
  \mathrm{(A)}\quad pH-m\rho\ge0\qquad\text{and}\qquad\mathrm{(B)}\quad pH+m\rho'\ge0,
\]
where $\rho=pt+2-p$ and $\rho'=2-p-pt$; note that $\rho+\rho'=2(2-p)>0$. The supporting line
at $s=\pm1$ and \cref{lem:ends} at $s=t$ give
\begin{align*}
  m(1-t)&\le h(1)-H,& \rho\,h(1)&\le2H,\\
  m(1+t)&\ge H-h(-1),& \rho'h(-1)&\le2H .
\end{align*}
If $\rho\ge0$, then $(1-t)(pH-m\rho)\ge pH(1-t)-\rho(h(1)-H)=2H-\rho\,h(1)\ge0$, which is
(A). If $\rho'\ge0$, then $(1+t)(pH+m\rho')\ge pH(1+t)+\rho'(H-h(-1))=2H-\rho'h(-1)\ge0$,
which is (B). Let $\rho<0$; then $\rho'>0$ and (B) holds. If $m\ge0$, then
$pH-m\rho\ge pH\ge0$; if $m<0$, then $pH-m\rho=(pH+m\rho')-m(\rho+\rho')\ge0$. This is (A).
The case $\rho'<0$ is symmetric.
\end{proof}

In terms of $a=(2-p)/p$, $\rho\ge0$ means $t\ge-a$ and $\rho'\ge0$ means $t\le a$.

\begin{proof}[Proof of \cref{lem:pone}]
By \cref{lem:ends} with $p=1$, $h(-1)=\tilde h(-1)\ge0$ and $h(1)\ge0$, and
$h(s)\ge\tfrac12h(1)(1+s)\ge0$ on $[-1,1]$.
\end{proof}

\begin{proof}[Proof of \cref{thm:low}]
By \cref{lem:profile} the profile $\Phi$ of $F$ is a nonnegative profile of degree $p$.
The Burkholder-type functions are positively $p$-homogeneous, so
\cref{prop:low} at $t=\det A_0/\abs{A_0}^2$ gives $c_0,m$ with
$G_{c_0,m}(\abs z,\abs w)\le\Phi(\abs z,\abs w)=F(A)$ for all $A$, with equality at $A_0$; for
$A_0=0$ take $c_0=m=0$. The formula for $G_{c_0,m}$ in terms of
$B_p(A)=\tfrac2p(\abs z-(p-1)\abs w)(\abs z+\abs w)^{p-1}$ and
$B_p(JA)=\tfrac2p(\abs w-(p-1)\abs z)(\abs z+\abs w)^{p-1}$ is an identity, and its coefficients
are nonnegative because $(2-p)\abs m\le pc_0$.
\end{proof}

\begin{proof}[Proof of \cref{cor:low}]
(a) The Burkholder function $B_p$ is quasiconvex, by \cite[Corollary~2]{SharpHI} as
formalized in~\cite{IwaniecLean}; see \cref{sec:intro-low} for the test maps. So is
$A\mapsto B_p(JA)$, by composing the test maps with $J$, and both are continuous. Hence every $G_{c_0,m}$ as in \cref{thm:low} is quasiconvex,
and for a test map $\varphi$,
\[
  F(A_0)=G_{c_0,m}(A_0)\le\int_{(0,1)^2}G_{c_0,m}(A_0+\nabla\varphi)\,dx
  \le\int_{(0,1)^2}F(A_0+\nabla\varphi)\,dx .
\]

(b) By \cref{lem:pone}, $F\ge0$, so \cref{thm:low} applies with $p=1$: at every
$A_0$ there are $c_0\ge\abs m$ such that $(c_0+m)\abs z+(c_0-m)\abs w$ lies below $F$ and equals
$F$ at $A_0$. This function is convex. A function with a convex touching minorant at every
point is convex.
\end{proof}

\subsection{The exponents \texorpdfstring{$0<p<1$}{0 < p < 1}}

\begin{proof}[Proof of \cref{lem:sublinear}]
Let $s<t$. For $T>t$, convexity gives
\[
  \frac{g(t)-g(s)}{t-s}\le\frac{g(T)-g(t)}{T-t},
\]
and the right side tends to $0$ as $T\to+\infty$, since $g(T)/T\to0$. For $T<s$,
\[
  \frac{g(s)-g(T)}{s-T}\le\frac{g(t)-g(s)}{t-s},
\]
and the left side tends to $0$ as $T\to-\infty$. Hence $g(s)=g(t)$.
\end{proof}

\begin{proof}[Proof of \cref{prop:small}]
A rank-one convex function is separately convex in the entries of the matrix, hence continuous.
From \ref{H2} with $\lambda=2$, $F(0)=2^pF(0)$, so $F(0)=0$.

Let $A\in\M$ and let $B$ have rank one, and put $g(\tau)=F(A+\tau B)$, which is convex. For
$\tau>0$, homogeneity gives $g(\tau)/\tau=\tau^{p-1}F(B+A/\tau)$, which tends to $0$ as
$\tau\to+\infty$ because $p<1$ and $F$ is continuous at $B$. The same with $-B$ in place of $B$
gives $g(\tau)/\tau\to0$ as $\tau\to-\infty$. By \cref{lem:sublinear}, $g$ is constant:
$F(A+\tau B)=F(A)$ for every $\tau$.

The four matrices with one entry equal to $1$ and the others $0$ have rank one. Removing the
entries of $A$ one at a time gives $F(A)=F(0)=0$.
\end{proof}

\begin{proof}[Proof of \cref{cor:all}]
For $0<p<1$, $F=0$ by \cref{prop:small}. For $p=1$, $F$ is convex by \cref{cor:low}\,(b), hence
quasiconvex. For $1<p<2$ this is
\cref{cor:low}\,(a). For $p\ge2$, $F$ is polyconvex by \cref{thm:main}, hence quasiconvex by
\cref{lem:supports}.
\end{proof}
